%% -*- coding: utf-8 -*-
\documentclass[12pt,pagesize,paper=192mm:108mm]{scrbook} 
%1920x1080 1280x720
\areaset[current]{192mm}{108mm}
\usepackage{calc}
\usepackage[T2A]{fontenc}
\usepackage[utf8]{inputenc}
\usepackage[english,russian]{babel}
\usepackage{microtype}
\usepackage{misccorr}
\usepackage{cmap}
%\usepackage[unicode=true]{hyperref}
\usepackage{graphicx}
\usepackage{amssymb}
\usepackage{amsmath}
%\usepackage{srcltx}
\usepackage{textcomp}
\usepackage{xspace}
%научные символы и смайлики \smiley \frownie
\usepackage{wasysym}
\usepackage{ccicons}
\begin{document}
\begin{titlepage}
  \vspace*{-0.5em}
  \begin{center}    
    \hspace*{3em}
    \begin{minipage}[t]{3em}
      \includegraphics[width=\textwidth]{../BINP-logo}
    \end{minipage}\hfill
    \begin{minipage}{0.23\linewidth}
    \includegraphics[width=\textwidth]{../NSU-logo}
    \end{minipage}
    \hfill
    \hspace*{6em}

    Кафедра теоретической физики физического факультета НГУ
    \medskip

    \Large
    Профессор Фадин В.\,С.
    \bigskip

    \huge
    \textbf{Квантовая электродинамика}
    \bigskip

    \Large
    Лекция № 9
    \vfill

    \normalsize
    % \begin{minipage}{0.65\linewidth}
    % \end{minipage}
    \vfill

    \normalsize \ccbysa\hspace{0.5em}  Новосибирск 2013
  \end{center}
\end{titlepage}
\vspace*{-1em}
\begin{center}
\vfill
  \begin{minipage}{0.9\linewidth}
    Асимптотика амплитуды в реджевском пределе из фейнмановских
    диаграмм: грибовское разложение метрического тензора для обмена
    фотоном в t"=канале, обмен скалярной и спинорной частицей. Процесс
    фотон-фотонного рассеяния: теорема Фарри, поведение амплитуды и
    сечения при малых энергиях, качественное поведение сечения с
    ростом энергии. Приближенные методы: факторизация сечения на
    жесткую (на малых пространственно"=временных масштабах) и мягкую
    часть (на больших пространственно"=временных масштабах). Рассеяние
    электрона во внешнем поле сопровождающееся излучением фотонов:
    матричный элемент без~излучения, матричный элемент для рассеяния с
    излучением фотона, мягкофотонное приближение, факторизация
    матричного элемента, фурье"=образ тока рассеивающегося электрона,
    сечение процесса рассеяния с излучением мягкого фотона,
    вероятность излучения одного фотона в борновском
    приближении. Обобщение на случай рассеяния произвольного числа
    начальных и конечных частиц. Обобщение на случай излучения
    нескольких фотонов. Связь с классическим выражением для
    спектральной плотности излучения. Проинтегрированная по углам
    вероятность излучения фотона (спектральная плотность):
    фейнмановское правило суммирования по поляризациям, удобные
    системы отсчета для вычисления интегралов по~углам, поведение в
    нерелятивистском и ультрарелятивистском пределах, существенная
    область углов излучения фотонов, роль интерференции между
    излучением фотона начальным и конечным электроном.
  \end{minipage}
  \vfill

  % \normalsize \ccbysa\hspace{0.5em} Новосибирск 2013
\end{center}
\end{document}
